% !TEX program = xelatex
% Assignment 5 - Implicit differentiation, valuation and first-order risk.
\documentclass[11pt]{article}
\usepackage[a4paper,top=1in,bottom=1in,left=1.1in,right=1.1in]{geometry}
\usepackage{setspace}
\setstretch{1.18}
\setlength{\parindent}{1.5em}
\setlength{\parskip}{2pt}
\setlength{\emergencystretch}{1.5em}
\usepackage{fontspec}
\setmainfont{texgyreheros-regular.otf}[BoldFont=texgyreheros-bold.otf,ItalicFont=texgyreheros-italic.otf,BoldItalicFont=texgyreheros-bolditalic.otf]
\setsansfont{texgyreheros-regular.otf}[BoldFont=texgyreheros-bold.otf]
\usepackage{amsmath,amssymb,booktabs,array,tabularx}
\usepackage{enumitem,xcolor,colortbl,titlesec,fancyhdr,float,caption}
\usepackage[hidelinks]{hyperref}
\definecolor{NYUviolet}{HTML}{57068C}
\definecolor{NYUviolet2}{HTML}{8900E1}
\definecolor{inkgrey}{HTML}{3A3A3A}
\definecolor{rulegrey}{HTML}{B9B9B9}
\definecolor{tablehead}{HTML}{ECE3F4}
\hypersetup{colorlinks=true,linkcolor=NYUviolet,urlcolor=NYUviolet2,citecolor=NYUviolet,
 pdftitle={Assignment 5: Implicit Differentiation and First-Order Risk},pdfauthor={Ziyi Yang}}
\titleformat{\section}{\sffamily\Large\bfseries\color{NYUviolet}}{\thesection}{0.7em}{}[\vspace{2pt}{\color{rulegrey}\titlerule[0.8pt]}]
\titleformat{\subsection}{\sffamily\large\bfseries\color{inkgrey}}{\thesubsection}{0.6em}{}
\titlespacing*{\section}{0pt}{16pt}{8pt}
\titlespacing*{\subsection}{0pt}{10pt}{5pt}
\pagestyle{fancy}
\fancyhf{}
\setlength{\headheight}{15pt}
\renewcommand{\headrulewidth}{0.6pt}
\fancyhead[L]{\footnotesize\sffamily\color{inkgrey}Implicit Differentiation and Risk}
\fancyhead[R]{\footnotesize\sffamily\color{inkgrey}FRE-GT-9743 -- Assignment 5}
\fancyfoot[C]{\footnotesize\sffamily\color{inkgrey}\thepage}
\setlist[itemize]{leftmargin=1.4em,itemsep=2pt,topsep=3pt,label=\textcolor{NYUviolet}{\textbullet}}
\setlist[enumerate]{leftmargin=1.6em,itemsep=2pt,topsep=3pt}
\newcolumntype{L}[1]{>{\raggedright\arraybackslash}p{#1}}
\newcolumntype{Y}{>{\raggedright\arraybackslash}X}
\renewcommand{\arraystretch}{1.12}
\newcommand{\keynote}[1]{\par\medskip\noindent\textbf{\sffamily\color{NYUviolet}Key result.}\enspace #1\par\medskip}
\newcommand{\focus}[1]{\noindent{\small\itshape\color{inkgrey}\textbf{Question focus.} #1}\par\medskip}
\numberwithin{equation}{section}
\captionsetup{font=small,labelfont={bf,sf,color=NYUviolet},labelsep=period}

\begin{document}
\setlength{\abovedisplayskip}{7pt}
\setlength{\belowdisplayskip}{7pt}
\setlength{\jot}{2pt}
\hypersetup{pageanchor=false}

\begin{titlepage}
\thispagestyle{empty}
\centering
\vspace*{0.55in}
{\color{NYUviolet}\rule{\textwidth}{2.2pt}}\\[10pt]
{\sffamily\fontsize{16}{19}\selectfont\color{inkgrey}Fixed Income Analysis Report}\\[14pt]
{\sffamily\bfseries\fontsize{25}{30}\selectfont\color{NYUviolet}
Implicit Differentiation and\\[3pt]First-Order Risk}\\[14pt]
{\sffamily\fontsize{15}{18}\selectfont\color{inkgrey}
Calibration Jacobians, Volatility Parameters,\\[3pt]
Discount-Factor Gradients and Cashflow Valuation}\\[10pt]
{\color{NYUviolet}\rule{\textwidth}{2.2pt}}\\[26pt]
\begin{minipage}{0.88\textwidth}
\centering\itshape\small\color{inkgrey}
This report solves the three written questions in Assignment 5 and documents
the completed valuation-engine implementation. It connects Newton calibration
to the implicit-function theorem, applies the framework to a three-parameter
implied-volatility smile, and verifies analytic curve-node risk for a fixed
accrued cashflow on a referenced instantaneous-forward-rate model.
\end{minipage}\\[34pt]
\begin{tabular}{rl}
\sffamily\bfseries Author & Ziyi Yang\\[3pt]
\sffamily\bfseries Course & FRE-GT-9743\\[3pt]
\sffamily\bfseries Source problem & Jianing (Jay) Yao, Assignment 5\\[3pt]
 & October 15, 2026
\end{tabular}
\vfill
{\sffamily\small\color{inkgrey}Assignment 5 \quad\textbullet\quad Calibration and Analytic Risk}\\[2pt]
{\color{NYUviolet}\rule{\textwidth}{1.2pt}}
\end{titlepage}
\hypersetup{pageanchor=true}

\begin{center}{\sffamily\large\bfseries\color{NYUviolet}Abstract}\end{center}
\begin{quote}\small\itshape
Calibration is a root-finding problem, while calibration risk is an implicit
differentiation problem. A nonsingular pricing Jacobian maps market-price
changes into model-parameter changes. For a volatility smile, this Jacobian is
the product of option vega and the interpolator's parameter gradient. In the
programming component, discount factors are exponentials of integrated
instantaneous forward rates. Their node gradients are propagated through a
referenced funding curve with the product rule and then scaled by the cashflow
amount to obtain exact first-order PV risk.
\end{quote}
\vspace{4pt}
{\small\noindent\textbf{\sffamily Keywords:}\;Implicit function \textbullet\; Newton calibration
\textbullet\; Volatility smile \textbullet\; IFR curve \textbullet\; First-order risk}
\vspace{8pt}{\color{rulegrey}\hrule}
\setcounter{tocdepth}{1}
\renewcommand{\contentsname}{\sffamily\large\bfseries\color{NYUviolet}Contents}
{\small\tableofcontents}
{\color{rulegrey}\hrule}
\clearpage

\section{Q1: Solving the Calibration Equations}
\focus{Give a general recipe for recovering the implied parameter vector $X^I$ from $V(X^I)=X^M$.}

Define the residual
\begin{equation}
R(X^I)=V(X^I)-X^M.
\end{equation}
Calibration seeks a root, $R(X^I)=0$. Starting from a reasonable guess
$X^I_0$, calculate the model Jacobian
\begin{equation}
J_V(X^I_n)=\left.\frac{\partial V}{\partial X^I}\right|_{X^I_n}
\end{equation}
and solve the Newton system
\begin{equation}
J_V(X^I_n)\,\Delta X^I_n=X^M-V(X^I_n).
\end{equation}
Update $X^I_{n+1}=X^I_n+\Delta X^I_n$ and repeat until both the pricing
residual and the parameter step satisfy the chosen tolerances. A damped step
or line search can be used when a full Newton step is unstable. Numerically,
the linear system should be solved directly rather than by forming an explicit
matrix inverse. Bounds or a reparameterization may be added when admissibility,
such as positive volatility, must be enforced.

Local uniqueness and a stable Newton step require a square, nonsingular
Jacobian near the solution. A poorly conditioned Jacobian signals that the
market observations do not identify the model parameters reliably.

\keynote{Calibration consists of defining price residuals, evaluating the pricing Jacobian, solving a linear Newton system, updating the parameters and checking convergence.}

\section{Q2: Sensitivity of Implied Parameters}
\focus{Calculate $dX^I/dX^M$ at the calibrated optimum.}

At the optimum, the calibration identity is
\begin{equation}
V(X^I(X^M))=X^M.
\end{equation}
Taking total differentials gives
\begin{equation}
J_V\,dX^I=dX^M.
\end{equation}
Provided $J_V$ is nonsingular, the implicit-function theorem yields
\begin{equation}
\boxed{\frac{dX^I}{dX^M}=J_V^{-1}.}
\end{equation}
Each column gives the response of all implied parameters to a unit change in
one market quote. In implementation, a particular shock is obtained by
solving $J_VdX^I=dX^M$ rather than explicitly inverting $J_V$.

For a downstream quantity $G(X^I)$, the chain rule gives market-quote risk:
\begin{equation}
\boxed{\frac{dG}{dX^M}=\frac{\partial G}{\partial X^I}J_V^{-1}.}
\end{equation}

\section{Q3: Three-Point Implied-Volatility Calibration}
\subsection{Calibration from option prices}
Let $B_i(\sigma)$ denote the undiscounted European option pricing function at
strike $K_i$. The three parameters $\theta\in\mathbb{R}^3$ are determined by
\begin{equation}
B_i\!\left(\sigma(K_i;\theta)\right)=P_i,
\qquad i=1,2,3.
\end{equation}
Equivalently, solve the three residual equations
\begin{equation}
R_i(\theta)=B_i\!\left(\sigma(K_i;\theta)\right)-P_i=0
\end{equation}
with the Newton procedure from Q1. One may also invert each option price to an
implied volatility first and then fit the three interpolator parameters; the
direct price-space formulation makes the market-price sensitivity immediate.

The calibration Jacobian is
\begin{equation}
J_{ij}=\frac{\partial P_i}{\partial\theta_j}
=\underbrace{\frac{\partial B_i}{\partial\sigma_i}}_{\mathrm{Vega}_i}
\frac{\partial\sigma(K_i;\theta)}{\partial\theta_j}.
\end{equation}
Therefore
\begin{equation}
\frac{d\theta}{dP}=J^{-1},
\qquad
\boxed{\frac{d\sigma(K)}{dP}
=\nabla_\theta\sigma(K;\theta)^\mathsf{T}J^{-1}.}
\end{equation}
The direct sensitivity of the interpolated volatility to its parameters is
simply $d\sigma(K)=\nabla_\theta\sigma(K;\theta)^\mathsf{T}d\theta$.

\subsection{Trader-interpretable parameters}
A natural three-parameter representation is level, skew and curvature:
\begin{align}
\theta_{\mathrm{ATM}}&=\sigma(K_{\mathrm{ATM}}),\\
\theta_{\mathrm{RR}}&=\sigma(K_3)-\sigma(K_1),\\
\theta_{\mathrm{BF}}&=\frac{\sigma(K_1)+\sigma(K_3)}{2}-\sigma(K_{\mathrm{ATM}}).
\end{align}
The first parameter moves the smile vertically, the risk reversal tilts the
two wings relative to one another, and the butterfly changes their average
height relative to ATM. This basis makes sensitivities easier to discuss than
sensitivities to three opaque polynomial coefficients.

\section{Completed Valuation Engine}
\subsection{Discount-factor gradient}
For one instantaneous-forward-rate component,
\begin{equation}
D(T)=\exp\left(-\int_0^T f(u)\,du\right).
\end{equation}
If $w_i(T)$ is the derivative of the integrated interpolant with respect to
node ordinate $f_i$, then
\begin{equation}
\boxed{\frac{\partial D(T)}{\partial f_i}=-D(T)w_i(T).}
\end{equation}
For the referenced funding curve
$D_{\mathrm{fund}}=D_{\mathrm{spread}}D_{\mathrm{SOFR}}$, the implementation
walks the reference chain and applies the product rule. Each component block
receives its local derivative multiplied by every other local discount factor.
The supplied \texttt{scaler} multiplies only the new contribution, while
\texttt{accumulate} controls replacement or addition.

\subsection{Fixed accrued cashflow}
For signed notional $N$ and accrual fraction $\alpha$, the amount is
$C=N\alpha$. A future payment has
\begin{equation}
PV=CD(T_p),\qquad \nabla PV=C\nabla D(T_p).
\end{equation}
A payment today has $PV=Cash=C$ and zero curve risk. A past payment has zero
PV, cash and risk.

\begin{table}[H]
\centering
\caption{Verified Assignment 5 outputs}
\small
\begin{tabularx}{\textwidth}{@{}L{4.1cm}Y@{}}
\toprule\rowcolor{tablehead}\textbf{Quantity}&\textbf{Result}\\\midrule
Cashflow PV & $245{,}870.910580$\\
Cash today & $0$\\
SOFR raw gradient & $(-60{,}625.703979,-123{,}272.264756,0)$\\
Funding-spread raw gradient & $(-121{,}925.026890,-61{,}972.941845)$\\
One-basis-point report & Raw gradients multiplied by $10^{-4}$, preserving sign\\
\bottomrule
\end{tabularx}
\end{table}
The zero in the third SOFR node is expected because its piecewise-constant
segment begins after the payment date.

\section{Conclusion}
The same Jacobian that drives Newton calibration also maps market shocks into
implied-parameter shocks. In the option example, vega connects price space to
volatility space and a level--skew--curvature parameterization makes the
resulting risk intuitive. The completed library applies the same derivative
discipline to yield curves: differentiate the integrated IFR, propagate the
gradient through referenced discount factors, and scale it by the contractual
cashflow. The executed notebook confirms the PV and every reported risk block.

\end{document}
